\chapter{WebDAV and Autoversioning}\label{svn.webdav}

WebDAV is an extension to HTTP, and it is growing more and more popular
as a standard for file sharing. Today\textquotesingle s operating
systems are becoming extremely web-aware, and many now have built-in
support for mounting ``shares'' exported by WebDAV servers.

If you use Apache as your Subversion network server, to some extent you
are also running a WebDAV server. This appendix gives some background on
the nature of this protocol, how Subversion uses it, and how well
Subversion interoperates with other software that is WebDAV-aware.

\section{What Is WebDAV?}\label{svn.webdav.basic}

DAV stands for ``Distributed Authoring and Versioning.'' RFC 2518
defines a set of concepts and accompanying extension methods to HTTP 1.1
that make the Web a more universal read/write medium. The basic idea is
that a WebDAV-compliant web server can act like a generic file server;
clients can ``mount'' shared folders over HTTP that behave much like
other network filesystems (such as NFS or SMB).

The tragedy, though, is that despite the acronym, the RFC specification
doesn\textquotesingle t actually describe any sort of version control.
Basic WebDAV clients and servers assume that only one version of each
file or directory exists, and that it can be repeatedly overwritten.

Because RFC 2518 left out versioning concepts, another committee was
left with the responsibility of writing RFC 3253 a few years later. The
new RFC adds versioning concepts to WebDAV, placing the ``V'' back in
``DAV''---hence the term ``DeltaV''. WebDAV/DeltaV clients and servers
are often called just ``DeltaV'' programs, since DeltaV implies the
existence of basic WebDAV.

The original WebDAV standard has been widely successful. Every modern
computer operating system has a general WebDAV client built in (details
to follow), and a number of popular standalone applications are also
able to speak WebDAV---Microsoft Office, Dreamweaver, and Photoshop, to
name a few. On the server end, Apache HTTP Server has been able to
provide WebDAV services since 1998 and is considered the de facto open
source standard. Several other commercial WebDAV servers are available,
including Microsoft\textquotesingle s own IIS.

DeltaV, unfortunately, has not been so successful. It\textquotesingle s
very difficult to find any DeltaV clients or servers. The few that do
exist are relatively unknown commercial products, and thus
it\textquotesingle s very difficult to test interoperability.
It\textquotesingle s not entirely clear as to why DeltaV has remained
stagnant. Some opine that the specification is just too complex. Others
argue that while WebDAV\textquotesingle s features have mass appeal
(even the least technical users appreciate network file sharing), its
version control features just aren\textquotesingle t interesting or
necessary for most users. Finally, some believe that DeltaV remains
unpopular because there\textquotesingle s still no open source server
product that implements it well.

When Subversion was still in its design phase, it seemed like a great
idea to use Apache as a network server. It already had a module to
provide WebDAV services. DeltaV was a relatively new specification. The
hope was that the Subversion server module (\texttt{mod\_dav\_svn})
would eventually evolve into an open source DeltaV reference
implementation. Unfortunately, DeltaV has a very specific versioning
model that doesn\textquotesingle t quite line up with
Subversion\textquotesingle s model. Some concepts were mappable; others
were not.

What does this mean, then?

First, the Subversion client is not a fully implemented DeltaV client.
It needs certain types of things from the server that DeltaV itself
cannot provide, and thus is largely dependent on a number of
Subversion-specific HTTP \texttt{REPORT} requests that only
\texttt{mod\_dav\_svn} understands.

Second, \texttt{mod\_dav\_svn} is not a fully realized DeltaV server.
Many portions of the DeltaV specification were irrelevant to Subversion,
and thus were left unimplemented.

A long-held debate in the Subversion developer community about whether
it was worthwhile to remedy either of these situations eventually
reached closure, with the Subversion developers officially deciding to
abandon plans to fully support DeltaV. As of Subversion 1.7, Subversion
clients and servers introduce numerous non-standard simplifications of
the DeltaV standards\footnote{The Subversion developers colloquially
  refer to this deviation from the DeltaV standard as ``HTTPv2''.}, with
more customizations of this sort likely to come. Those versions of
Subversion will, of course, continue to provide the same DeltaV feature
support already present in older releases, but no new work will be done
to increase coverage of the specification---Subversion is intentionally
moving away from strict DeltaV as its primary HTTP-based protocol.

\section{Autoversioning}\label{svn.webdav.autoversioning}

While the Subversion client is not a full DeltaV client, and the
Subversion server is not a full DeltaV server, there\textquotesingle s
still a glimmer of WebDAV interoperability to be happy about:
autoversioning.

Autoversioning is an optional feature defined in the DeltaV standard. A
typical DeltaV server will reject an ignorant WebDAV client attempting
to do a \texttt{PUT} to a file that\textquotesingle s under version
control. To change a version-controlled file, the server expects a
series of proper versioning requests: something like
\texttt{MKACTIVITY}, \texttt{CHECKOUT}, \texttt{PUT}, \texttt{CHECKIN}.
But if the DeltaV server supports autoversioning, write requests from
basic WebDAV clients are accepted. The server behaves as though the
client \emph{had} issued the proper series of versioning requests,
performing a commit under the hood. In other words, it allows a DeltaV
server to interoperate with ordinary WebDAV clients that
don\textquotesingle t understand versioning.

Because so many operating systems already have integrated WebDAV
clients, the use case for this feature can be incredibly appealing to
administrators working with non-technical users. Imagine an office of
ordinary users running Microsoft Windows or Mac OS. Each user ``mounts''
the Subversion repository, which appears to be an ordinary network
folder. They use the shared folder as they always do: open files, edit
them, and save them. Meanwhile, the server is automatically versioning
everything. Any administrator (or knowledgeable user) can still use a
Subversion client to search history and retrieve older versions of data.

This scenario isn\textquotesingle t fiction---it\textquotesingle s real
and it works. To activate autoversioning in \texttt{mod\_dav\_svn}, use
the \texttt{SVNAutoversioning} directive within the \texttt{httpd.conf}
\texttt{Location} block, like so:

\begin{verbatim}
<Location /repos>
  DAV svn
  SVNPath /var/svn/repository
  SVNAutoversioning on
</Location>
\end{verbatim}

When Subversion autoversioning is active, write requests from WebDAV
clients result in automatic commits. A generic log message is
automatically generated and attached to each revision.

Before activating this feature, however, understand what
you\textquotesingle re getting into. WebDAV clients tend to do
\emph{many} write requests, resulting in a huge number of automatically
committed revisions. For example, when saving data, many clients will do
a \texttt{PUT} of a 0-byte file (as a way of reserving a name) followed
by another \texttt{PUT} with the real file data. The single file-write
results in two separate commits. Also consider that many applications
auto-save every few minutes, resulting in even more commits.

If you have a post-commit hook program that sends email, you may want to
disable email generation either altogether or on certain sections of the
repository; it depends on whether you think the influx of emails will
still prove to be valuable notifications or not. Also, a smart
post-commit hook program can distinguish between a transaction created
via autoversioning and one created through a normal Subversion commit
operation. The trick is to look for a revision property named
\texttt{svn:autoversioned}. If present, the commit was made by a generic
WebDAV client.

Another feature that may be a useful complement for
Subversion\textquotesingle s autoversioning comes from
Apache\textquotesingle s \texttt{mod\_mime} module. If a WebDAV client
adds a new file to the repository, there\textquotesingle s no
opportunity for the user to set the the \texttt{svn:mime-type} property.
This might cause the file to appear as a generic icon when viewed within
a WebDAV shared folder, not having an association with any application.
One remedy is to have a sysadmin (or other Subversion-knowledgeable
person) check out a working copy and manually set the
\texttt{svn:mime-type} property on necessary files. But
there\textquotesingle s potentially no end to such cleanup tasks.
Instead, you can use the \texttt{ModMimeUsePathInfo} directive in your
Subversion \texttt{\textless{}Location\textgreater{}} block:

\begin{verbatim}
<Location /repos>
  DAV svn
  SVNPath /var/svn/repository
  SVNAutoversioning on

  ModMimeUsePathInfo on

</Location>
\end{verbatim}

This directive allows \texttt{mod\_mime} to attempt automatic deduction
of the MIME type on new files that enter the repository via
autoversioning. The module looks at the file\textquotesingle s named
extension and possibly the contents as well; if the file matches some
common patterns, the file\textquotesingle s \texttt{svn:mime-type}
property will be set automatically.

\section{Client Interoperability}\label{svn.webdav.clients}

All WebDAV clients fall into one of three categories---standalone
applications, file-explorer extensions, or filesystem implementations.
These categories broadly define the types of WebDAV functionality
available to users. \hyperref[svn.webdav.clients.tbl-1]{Common WebDAV
clients} gives our categorization as well as a quick description of some
common pieces of WebDAV-enabled software. You can find more details
about these software offerings, as well as their general category, in
the sections that follow.

\begin{longtable}[]{@{}>{\raggedright\arraybackslash}p{\dimexpr(\linewidth-2\tabcolsep*6)/6\relax}>{\raggedright\arraybackslash}p{\dimexpr(\linewidth-2\tabcolsep*6)/6\relax}>{\raggedright\arraybackslash}p{\dimexpr(\linewidth-2\tabcolsep*6)/6\relax}>{\raggedright\arraybackslash}p{\dimexpr(\linewidth-2\tabcolsep*6)/6\relax}>{\raggedright\arraybackslash}p{\dimexpr(\linewidth-2\tabcolsep*6)/6\relax}>{\raggedright\arraybackslash}p{\dimexpr(\linewidth-2\tabcolsep*6)/6\relax}@{}}
\caption{Common WebDAV
clients}\label{svn.webdav.clients.tbl-1}\tabularnewline
\toprule\noalign{}
Software & Type & Windows & Mac & Linux & Description \\
\midrule\noalign{}
\endfirsthead
\toprule\noalign{}
Software & Type & Windows & Mac & Linux & Description \\
\midrule\noalign{}
\endhead
\bottomrule\noalign{}
\endlastfoot
Adobe Photoshop & Standalone WebDAV application & X & & & Image editing
software, allowing direct opening from, and writing to, WebDAV URLs \\
cadaver & Standalone WebDAV application & & X & X & Command-line WebDAV
client supporting file transfer, tree, and locking operations \\
DAV Explorer & Standalone WebDAV application & X & X & X & Java GUI tool
for exploring WebDAV shares \\
Adobe Dreamweaver & Standalone WebDAV application & X & & & Web
production software able to directly read from and write to WebDAV
URLs \\
Microsoft Office & Standalone WebDAV application & X & & & Office
productivity suite with several components able to directly read from
and write to WebDAV URLs \\
Microsoft Web Folders & File-explorer WebDAV extension & X & & & GUI
file explorer program able to perform tree operations on a WebDAV
share \\
GNOME Nautilus & File-explorer WebDAV extension & & & X & GUI file
explorer able to perform tree operations on a WebDAV share \\
KDE Konqueror & File-explorer WebDAV extension & & & X & GUI file
explorer able to perform tree operations on a WebDAV share \\
Mac OS X & WebDAV filesystem implementation & & X & & Operating system
that has built-in support for mounting WebDAV shares. \\
Novell NetDrive & WebDAV filesystem implementation & X & & &
Drive-mapping program for assigning Windows drive letters to a mounted
remote WebDAV share \\
SRT WebDrive & WebDAV filesystem implementation & X & & & File transfer
software, which, among other things, allows the assignment of Windows
drive letters to a mounted remote WebDAV share \\
davfs2 & WebDAV filesystem implementation & & & X & Linux filesystem
driver that allows you to mount a WebDAV share \\
\end{longtable}

\subsection{Standalone WebDAV
Applications}\label{svn.webdav.clients.standalone}

A WebDAV application is a program that speaks WebDAV protocols with a
WebDAV server. We\textquotesingle ll cover some of the most popular
programs with this kind of WebDAV support.

\subsubsection{Microsoft Office, Dreamweaver,
Photoshop}\label{svn.webdav.clients.standalone.windows}

On Windows, several well-known applications contain integrated WebDAV
client functionality, such as Microsoft\textquotesingle s
Office,\footnote{WebDAV support was removed from Microsoft Access for
  some reason, but it exists in the rest of the Office suite.}
Adobe\textquotesingle s Photoshop and Dreamweaver programs.
They\textquotesingle re able to directly open and save to URLs, and tend
to make heavy use of WebDAV locks when editing a file.

Note that while many of these programs also exist for Mac OS X, they do
not appear to support WebDAV directly on that platform. In fact, on Mac
OS X, the File→Open dialog box doesn\textquotesingle t allow one to type
a path or URL at all. It\textquotesingle s likely that the WebDAV
features were deliberately left out of Macintosh versions of these
programs, since OS X already provides such excellent low-level
filesystem support for WebDAV.

\subsubsection{cadaver, DAV
Explorer}\label{svn.webdav.clients.standalone.free}

cadaver is a bare-bones Unix command-line program for browsing and
changing WebDAV shares. It uses the neon HTTP library---not
surprisingly, since both neon and cadaver are written by the same
author. cadaver is free software (GPL license) and is available at
\href{https://notroj.github.io/cadaver/}{}.

Using cadaver is similar to using a command-line FTP program, and thus
it\textquotesingle s extremely useful for basic WebDAV debugging. It can
be used to upload or download files in a pinch, to examine properties,
and to copy, move, lock, or unlock files:

\begin{verbatim}
$ cadaver http://host/repos
dav:/repos/> ls
Listing collection `/repos/': succeeded.
Coll: > foobar                                 0  May 10 16:19
      > playwright.el                       2864  May  4 16:18
      > proofbypoem.txt                     1461  May  5 15:09
      > westcoast.jpg                      66737  May  5 15:09

dav:/repos/> put README
Uploading README to `/repos/README':
Progress: [=============================>] 100.0% of 357 bytes succeeded.

dav:/repos/> get proofbypoem.txt
Downloading `/repos/proofbypoem.txt' to proofbypoem.txt:
Progress: [=============================>] 100.0% of 1461 bytes succeeded.
\end{verbatim}

DAV Explorer is another standalone WebDAV client, written in Java.
It\textquotesingle s under a free Apache-like license and is available
at \href{https://www.davexplorer.org/}{}. It does everything cadaver
does, but has the advantages of being portable and being a more
user-friendly GUI application. It\textquotesingle s also one of the
first clients to support the new WebDAV Access Control Protocol (RFC
3744).

Of course, DAV Explorer\textquotesingle s ACL support is useless in this
case, since \texttt{mod\_dav\_svn} doesn\textquotesingle t support it.
The fact that both cadaver and DAV Explorer support some limited DeltaV
commands isn\textquotesingle t particularly useful either, since they
don\textquotesingle t allow \texttt{MKACTIVITY} requests. But
it\textquotesingle s not relevant anyway; we\textquotesingle re assuming
all of these clients are operating against an autoversioning repository.

\subsection{File-Explorer WebDAV
Extensions}\label{svn.webdav.clients.file-explorer-extensions}

Some popular file explorer GUI programs support WebDAV extensions that
allow a user to browse a DAV share as though it was just another
directory on the local computer, and to perform basic tree editing
operations on the items in that share. For example, Windows Explorer is
able to browse a WebDAV server as a ``network place.'' Users can drag
files to and from the desktop, or can rename, copy, or delete files in
the usual way. But because it\textquotesingle s only a feature of the
file explorer, the DAV share isn\textquotesingle t visible to ordinary
applications. All DAV interaction must happen through the explorer
interface.

\subsubsection{Microsoft Web
Folders}\label{svn.webdav.clients.file-explorer-extensions.windows}

Microsoft was one of the original backers of the WebDAV specification,
and first started shipping a client in Windows 98, which was known as
Web Folders. This client was also shipped in Windows NT 4.0 and Windows
2000.

The original Web Folders client was an extension to Explorer, the main
GUI program used to browse filesystems. It works well enough. In Windows
98, the feature might need to be explicitly installed if Web Folders
aren\textquotesingle t already visible inside My Computer. In Windows
2000, simply add a new ``network place,'' enter the URL, and the WebDAV
share will pop up for browsing.

With the release of Windows XP, Microsoft started shipping a new
implementation of Web Folders, known as the WebDAV Mini-Redirector. The
new implementation is a filesystem-level client, allowing WebDAV shares
to be mounted as drive letters. Unfortunately, this implementation is
incredibly buggy. The client usually tries to convert HTTP URLs
(\texttt{http://host/repos}) into UNC share notation
(\texttt{\textbackslash{}\textbackslash{}host\textbackslash{}repos}); it
also often tries to use Windows Domain authentication to respond to
basic-auth HTTP challenges, sending usernames as
\texttt{HOST\textbackslash{}username}. These interoperability problems
are severe and are documented in numerous places around the Web, to the
frustration of many users. Even Greg Stein, the original author of
Apache\textquotesingle s WebDAV module, bluntly states that XP Web
Folders simply can\textquotesingle t operate against an Apache server.

Windows Vista\textquotesingle s initial implementation of Web Folders
seems to be almost the same as XP\textquotesingle s, so it has the same
sort of problems. With luck, Microsoft will remedy these issues in a
Vista Service Pack.

However, there seem to be workarounds for both XP and Vista that allow
Web Folders to work against Apache. Users have mostly reported success
with these techniques, so we\textquotesingle ll relay them here.

On Windows XP, you have two options. First, search
Microsoft\textquotesingle s web site for update KB907306, ``Software
Update for Web Folders.'' This may fix all your problems. If it
doesn\textquotesingle t, it seems that the original pre-XP Web Folders
implementation is still buried within the system. You can unearth it by
going to Network Places and adding a new network place. When prompted,
enter the URL of the repository, but \emph{include a port number} in the
URL. For example, you should enter \texttt{http://host/repos} as
\texttt{http://host:80/repos} instead. Respond to any authentication
prompts with your Subversion credentials.

On Windows Vista, the same KB907306 update may clear everything up. But
there may still be other issues. Some users have reported that Vista
considers all \texttt{http://} connections insecure, and thus will
always fail any authentication challenges from Apache unless the
connection happens over \texttt{https://}. If you\textquotesingle re
unable to connect to the Subversion repository via SSL, you can tweak
the system registry to turn off this behavior. Just change the value of
the
\texttt{HKEY\_LOCAL\_MACHINE\textbackslash{}SYSTEM\textbackslash{}CurrentControlSet\textbackslash{}Services\textbackslash{}WebClient\textbackslash{}Parameters\textbackslash{}BasicAuthLevel}
key from \texttt{1} to \texttt{2}. A final warning: be sure to set up
the Web Folders to point to the repository\textquotesingle s root
directory (\texttt{/}), rather than some subdirectory such as
\texttt{/trunk}. Vista Web Folders seems to work only against repository
roots.

In general, while these workarounds may function for you, you might get
a better overall experience using a third-party WebDAV client such as
WebDrive or NetDrive.

\subsubsection{Nautilus,
Konqueror}\label{svn.webdav.clients.file-explorer-extensions.linux-de}

Nautilus is the official file manager/browser for the GNOME desktop
(\href{https://www.gnome.org}{}), and Konqueror is the manager/browser
for the KDE desktop (\href{https://kde.org}{}). Both of these
applications have an explorer-level WebDAV client built in, and they
operate just fine against an autoversioning repository.

In GNOME\textquotesingle s Nautilus, select the File→Open location menu
item and enter the URL in the dialog box presented. The repository
should then be displayed like any other filesystem.

In KDE\textquotesingle s Konqueror, you need to use the
\texttt{webdav://} scheme when entering the URL in the location bar. If
you enter an \texttt{http://} URL, Konqueror will behave like an
ordinary web browser. You\textquotesingle ll likely see the generic HTML
directory listing produced by \texttt{mod\_dav\_svn}. When you enter
\texttt{webdav://host/repos} instead of \texttt{http://host/repos},
Konqueror becomes a WebDAV client and displays the repository as a
filesystem.

\subsection{WebDAV Filesystem
Implementation}\label{svn.webdav.clients.fs-impl}

The WebDAV filesystem implementation is arguably the best sort of WebDAV
client. It\textquotesingle s implemented as a low-level filesystem
module, typically within the operating system\textquotesingle s kernel.
This means that the DAV share is mounted like any other network
filesystem, similar to mounting an NFS share on Unix or attaching an SMB
share as a drive letter in Windows. As a result, this sort of client
provides completely transparent read/write WebDAV access to all
programs. Applications aren\textquotesingle t even aware that WebDAV
requests are happening.

\subsubsection{WebDrive,
NetDrive}\label{svn.webdav.clients.fs-impl.windows}

Both WebDrive and NetDrive are excellent commercial products that allow
a WebDAV share to be attached as drive letters in Windows. As a result,
you can operate on the contents of these WebDAV-backed pseudodrives as
easily as you can against real local hard drives, and in the same ways.
You can purchase WebDrive from South River Technologies
(\href{https://webdrive.com/}{}). Novell\textquotesingle s NetDrive is
freely available online, but requires users to have a NetWare license.

\subsubsection{Mac OS X}\label{svn.webdav.clients.fs-impl.macosx}

Apple\textquotesingle s OS X operating system has an integrated
filesystem-level WebDAV client. From the Finder, select the Go→Connect
to Server menu item. Enter a WebDAV URL, and it appears as a disk on the
desktop, just like any other mounted volume. You can also mount a WebDAV
share from the Darwin terminal by using the \texttt{webdav} filesystem
type with the \texttt{mount} command:

\begin{verbatim}
$ mount -t webdav http://svn.example.com/repos/project /some/mountpoint
$
\end{verbatim}

Note that if your \texttt{mod\_dav\_svn} is older than version 1.2, OS X
will refuse to mount the share as read/write; it will appear as
read-only. This is because OS X insists on locking support for
read/write shares, and the ability to lock files first appeared in
Subversion 1.2.

Also, OS X\textquotesingle s WebDAV client can sometimes be overly
sensitive to HTTP redirects. If OS X is unable to mount the repository
at all, you may need to enable the \texttt{BrowserMatch} directive in
the Apache server\textquotesingle s \texttt{httpd.conf}:

\begin{verbatim}
BrowserMatch "^WebDAVFS/1.[012]" redirect-carefully
\end{verbatim}

\subsubsection{Linux davfs2}\label{svn.webdav.clients.fs-impl.linux}

Linux davfs2 is a filesystem module for the Linux kernel, whose
development is organized at
\href{https://savannah.nongnu.org/projects/davfs2}{}. Once you install
davfs2, you can mount a WebDAV network share using the usual Linux mount
command:

\begin{verbatim}
$ mount.davfs http://host/repos /mnt/dav
\end{verbatim}

